\chapter*{附录三 康熙推迟重立太子谕}\addcontentsline{toc}{chapter}{附录三 康熙推迟重立太子谕}\markboth{附录三 康熙推迟重立太子谕}{附录三 康熙推迟重立太子谕}
% 原书第381页

赵申乔奏：皇太子乃国本，务宜早立。


% 原书第381页

朕自幼读书，诸凡留心，岂有立皇太子之要事而忘之理乎？因事关重大，未易轻立。曩者立太子时，尀耐索额图蓄积私心，倡言议定，以皇太子所制用，皆如朕制，俨若二君。天无二日，民无二王，其所以骄矜者，即此故耳。索额图者，我家之罪魁也。《大清会典》所载皇太子仪制，照古酌量定之可也。


% 原书第381–382页

又曰：宋仁宗三十年来，未立太子。我太祖皇帝未立太子，诸王贝勒大臣请太宗皇帝登极。太宗皇帝并未以世祖立为太子。汉唐以来，太子幼小，在祚不甚久者，犹且小事。若太子长大，下人结成党类，久而未有不失事者。人非圣贤，不能无过。尽其子职，如武王之子，安可得欤！至于众阿哥们，其所学所识，不为不及，俱已长大，各自分居。无奈伊等下人莫不各为其主，即册立，焉能保其事。此福不易享受。伊等亦必有企望之心，即此非宜。况为君难，为臣不易。常观古来为君之人，穷兵黩武者有之，好佛老者有之，尚刑罚者有之。朕惟宵旰恐惧，上思祖宗垂统之重，下体臣民仰望之心，体恤元元，五十余年，幸无外患敌国，亦可谓之能守矣。今若册立太子，必如克副朕之此心者，始可为之，岂宜造次？


% 原书第382页

且太子系为国本，国本不正，关系非浅。皇太子朕自幼教训，并不挞人骂人，嗣后长大，变成残暴，行事乖张，曷可胜言。详审其由，亦因病之所魔，至于昏聩耳。自废黜以至复立，此六年以来，因为教伊，朕之心血尽矣，须发皓矣。朕始终惟望其痊愈耳，并非别有册立之心而废黜。果有此心，岂可废而复立？父子之间不责善，责善则离，离则不祥莫大焉！又云：人莫知其子之恶，莫知其苗之硕。溺爱不明，贪得者无厌。处父子之间，严固非宜，宽亦不可，实系难事。朕惟思能得善处之宜。如卿等亦各有子。凡人之子，自幼教训，至于长大，各从朋类，别有交游，有不便于拘管者。


% 原书第382–383页

近如赵申乔，居官尚为清廉，朕问伊子太原知府赵凤诏做官如何。彼奏云：做官还好。后来闻得赵凤诏做官甚属平常。据此看来，人不能不存一点私心。因将赵申乔所奏折子，交与大学士温达去。册立太子之事，未可轻定。业已立错，再若轻定，复至于错，为之奈何？朕欲此意晓谕卿等，故聚集之，将此折交还他。


% 原书第383页

以上文字出自《本朝题驳公案》，系当时官场邸报所刊服本。相似内容另见《圣祖仁皇帝实录》卷253，康熙五十二年二月庚戌。

